The exercises below come from the ten tutorial series accompanying the course (by A.~Pilitowska and
T.~Brengos), collected here in the notation of this book; the three complex-numbers series are merged
into a single section. Exercises marked with ($\ast$) are more challenging.

\clearpage
\section{Operations modulo \texorpdfstring{$n$}{n}. Groups and fields}

\begin{exercise}
Let $n$ be a fixed positive integer and let $a$ and $b$ be any two integers. Recall that $a$ is congruent
to $b$ modulo $n$ if $n$ divides $a-b$. Calculate:
\begin{enumerate*}[label=\alph*), itemjoin={\qquad}]
\item $17 \bmod 4$
\item $4 \bmod 17$
\item $(-2) \bmod 5$
\end{enumerate*}
\end{exercise}

\begin{exercise}
Solve the equations:
\begin{enumerate}[label=\alph*)]
\item $2x \equiv 1 \pmod 7$
\item $2x \equiv 1 \pmod 6$
\item $5x \equiv 1 \pmod 9$
\item $x^2 \equiv 3 \pmod{11}$
\item $2x+3 \equiv 0 \pmod 5$
\item $x+k \equiv 0 \pmod n$
\end{enumerate}
\end{exercise}

\begin{exercise}[$\ast$]
Prove that for all $n,k,p\in \Z$ we have the following:
\begin{enumerate}[label=\alph*)]
\item $(n \bmod k) \bmod k = n \bmod k$
\item $(n+p) \bmod k = \big((n \bmod k) + (p \bmod k)\big) \bmod k$
\item $(np) \bmod k = \big((n \bmod k)\,(p \bmod k)\big) \bmod k$
\end{enumerate}
\end{exercise}

\begin{exercise}
Recall that for $a,b\in\N$ the operations of addition and multiplication modulo $n$ are defined by
$a +_n b = (a+b) \bmod n$ and $a \cdot_n b = (a\cdot b) \bmod n$. Prove that:
\begin{enumerate}[label=\alph*)]
\item addition modulo $n$ is associative, i.e.\ $(a +_n b) +_n c = a +_n (b +_n c)$ for any $a,b,c\in\N$;
\item multiplication modulo $n$ is associative, i.e.\ $(a \cdot_n b) \cdot_n c = a \cdot_n (b \cdot_n c)$
for any $a,b,c\in\N$;
\item multiplication modulo $n$ is distributive with respect to addition modulo $n$, i.e.\
$(a +_n b) \cdot_n c = (a \cdot_n c) +_n (b \cdot_n c)$ for any $a,b,c\in\N$.
\end{enumerate}
\end{exercise}

\begin{exercise}
State which of the following are groups:
\begin{enumerate*}[label=\alph*), itemjoin={\qquad}]
\item $(\Z_n, +_n)$
\item $(\Z_n\setminus\{0\}, \cdot_n)$
\item $(\{1,2,3,4,6,12\}, \gcd)$
\end{enumerate*}
\end{exercise}

\begin{exercise}[$\ast$]
Show that $(2^X, \bigtriangleup)$ is a group, where $\bigtriangleup$ denotes the symmetric difference.
Construct the group table of $(2^{\{a,b,c\}}, \bigtriangleup)$.
\end{exercise}

\begin{exercise}
Show that $(\Z_n\setminus\{0\}, \cdot_n)$ is a group if and only if $n$ is a prime number.
\end{exercise}

\begin{exercise}
Verify whether $(\R_+, \bullet)$ is a group, where $a\bullet b = a^2 b^2$.
\end{exercise}

\begin{exercise}
Let $p\in\N$ be a prime number. Show that $(\Z_p, +_p, \cdot_p)$ is a field.
\end{exercise}

\begin{exercise}
Determine which of the following structures are fields. Do this in two steps: first verify whether the
set is a group with respect to the first operation, and then whether the set without the identity element
of the first operation is a group with respect to the second operation. Here $2^X$ carries the operations
$\cup$, $\cap$ and the symmetric difference $\bigtriangleup$, while $\R^\R$ denotes the set of all
functions $\R\to\R$, with pointwise addition $+$ and composition $\circ$.
\begin{enumerate*}[label=\alph*), itemjoin={\qquad}]
\item $(2^X, \cup, \cap)$
\item $(\R^\R, +, \circ)$
\item $(2^X, \cap, \cup)$
\item $(\R^\R, \circ, +)$
\item $(2^X, \cap, \bigtriangleup)$
\item $(2^X, \bigtriangleup, \cap)$
\item $(2^X, \cup, \bigtriangleup)$
\item $(2^X, \bigtriangleup, \cup)$
\end{enumerate*}
\end{exercise}

\begin{exercise}
Verify whether $(\R, \oplus, \odot)$ is a field, where $a \oplus b = a+b+1$ and $a\odot b = ab+a+b$.
\end{exercise}

\clearpage
\section{Complex numbers}

\begin{exercise}
Calculate:
\begin{enumerate*}[label=\alph*), itemjoin={\qquad}]
\item $(2+3i)(4-2i)$
\item $(1+2i)(2i-1)$
\item $(3+i)(3-i)$
\item $(a+bi)(a-bi)$
\end{enumerate*}
\end{exercise}

\begin{exercise}
Calculate:
\begin{enumerate*}[label=\alph*), itemjoin={\qquad}]
\item $\dfrac{2-3i}{3+4i}$
\item $\dfrac{i}{1+i}$
\item $\dfrac{3-4i}{i}$
\end{enumerate*}
\end{exercise}

\begin{exercise}
Calculate:
\begin{enumerate*}[label=\alph*), itemjoin={\qquad}]
\item $\dfrac{1+4i}{4-i}$
\item $\dfrac{(3+4i)^4}{(3-4i)^3}$
\end{enumerate*}
\end{exercise}

\begin{exercise}
Show that for every two complex numbers $z$ and $w$:
\begin{enumerate}[label=\alph*)]
\item $\overline{z+w} = \bar z + \bar w$ and $\overline{zw} = \bar z\, \bar w$
\item $|zw| = |z|\,|w|$
\item $\left|\dfrac{z}{w}\right| = \dfrac{|z|}{|w|}$
\end{enumerate}
\end{exercise}

\begin{exercise}[$\ast$]
Show that multiplication of complex numbers is associative.
\end{exercise}

\begin{exercise}[$\ast$]
Show that multiplication of complex numbers is distributive with respect to addition.
\end{exercise}

\begin{exercise}
Describe geometrically the sets of points $z\in\C$ satisfying the following conditions:
\begin{multicols}{2}
\begin{enumerate}[label=\alph*)]
\item $|z-i+3| = 5$
\item $|z-i+3| > 5$
\item $|z-1+3i| \leq 5$
\item $2\leq |z+i| < 4$
\item $\Re\big((z-i)^2\big) = 0$
\item $\Re(z^2) = 2$ and $\big(\Im(z+i)\big)^2 = 1$
\item $\left|\dfrac{z+3}{z-2i}\right| \geq 1$
\item $\Re(z+1)<0$ and $|i-z|\leq 3$
\end{enumerate}
\end{multicols}
\end{exercise}

\begin{exercise}
State which of the following are groups:
\begin{enumerate*}[label=\alph*), itemjoin={\qquad}]
\item $(\{z\in\C \mid |z|=1\}, +)$
\item $(\{z\in\C \mid |z|=1\}, \cdot)$
\end{enumerate*}
\end{exercise}

\begin{exercise}
Find polar forms of the following complex numbers:
\begin{multicols}{2}
\begin{enumerate}[label=\alph*)]
\item $1+i$
\item $-3$ (negative $3$)
\item $-i$
\item $2+2i$
\item $3+3\sqrt3\, i$
\item $z = \sin\alpha + i\cos\alpha$
\item $1 - i\sqrt3$
\item $-z$, if $z = \cos\alpha + i \sin\alpha$
\item $\sqrt3 - i$
\item $\bar z$, if $z = \cos\alpha + i \sin\alpha$
\item $\tfrac12 + \tfrac{\sqrt3}{2}\, i$
\end{enumerate}
\end{multicols}
\end{exercise}

\begin{exercise}
Calculate:
\begin{multicols}{2}
\begin{enumerate}[label=\alph*)]
\item $\sqrt[3]{(1+i)^6}$
\item $\left(\dfrac{1-i}{1+i}\right)^{222}$
\item $\left(\dfrac12 + i\,\dfrac{\sqrt3}{2}\right)^{124}$
\item $\left(\dfrac{2+2i}{\sqrt3 - i}\right)^{8}$
\item $\sqrt[6]{1}$
\item $\sqrt[4]{-1}$
\item $\sqrt[4]{(3-5i)^4}$
\item $\sqrt[6]{-1}$
\end{enumerate}
\end{multicols}
\end{exercise}

\begin{exercise}
Describe geometrically the sets of all the following roots:
\begin{enumerate*}[label=\alph*), itemjoin={\qquad}]
\item $\sqrt[6]{1}$
\item $\sqrt[4]{-1}$
\item $\sqrt[8]{1}$
\item $\sqrt[6]{-1}$
\end{enumerate*}
\end{exercise}

\begin{exercise}
Show that for every $n$ all complex roots of $1$ of order $n$ form a group under multiplication. Write
down the table of the group of all complex roots of $1$ of order $4$.
\end{exercise}

\begin{exercise}
Use De Moivre's formula to:
\begin{enumerate}[label=(\alph*)]
\item express $\sin 5\alpha$ in terms of $\sin\alpha$ and $\cos\alpha$,
\item express $\cos 4\alpha$ in terms of $\sin\alpha$ and $\cos\alpha$.
\end{enumerate}
\end{exercise}

\begin{exercise}[$\ast$]
Prove that
\[
\frac{\cos\alpha + i \sin\alpha}{\cos\beta + i\sin\beta} = \cos(\alpha-\beta) + i \sin(\alpha - \beta).
\]
\end{exercise}

\begin{exercise}
Find the polar form of:
\begin{enumerate*}[label=\alph*), itemjoin={\qquad}]
\item $2+\sqrt2 + i \sqrt 2$
\item $2 + i + i\sqrt3$
\end{enumerate*}
\end{exercise}

\begin{exercise}
Calculate:
\begin{enumerate*}[label=\alph*), itemjoin={\qquad}]
\item $\cos\tfrac{\pi}{12}$
\item $\cos 105^\circ$
\item $\cos 75^\circ$
\item $\cos(-15^\circ)$
\end{enumerate*}
\end{exercise}

\begin{exercise}[$\ast$]
Calculate $(1+\cos\alpha+i\sin\alpha)^{10}$ for $\alpha\in [0,\pi)$.
\end{exercise}

\begin{exercise}
Let $f(x)\in\R[x]$. Show that $z$ is a root of $f$ if and only if $\bar z$ is a root of $f$.
\end{exercise}

\begin{exercise}
Factor each of the following polynomials into a product of factors of degree at most $2$:
\begin{enumerate*}[label=\alph*), itemjoin={\qquad}]
\item $x^6+1$
\item $x^4+1$
\item $x^6+x^3+1$
\item $x^5+1$
\end{enumerate*}
\end{exercise}

\begin{exercise}
Prove that every polynomial of an odd degree, with all real coefficients, has at least one real root.
\end{exercise}

\begin{exercise}
Knowing that $1-2i$ is a root of the polynomial $z^4+2z^3+2z^2+10z+25$, find all remaining roots.
\end{exercise}

\begin{exercise}
Knowing that $2+i$ is a root of the polynomial $z^4-2z^3-6z^2+22z-15$, find all remaining roots.
\end{exercise}

\begin{exercise}
Solve the equations:
\begin{enumerate*}[label=\alph*), itemjoin={\qquad}]
\item $x^2+x+1 = 0$
\item $2x^2+ix+3 = 0$
\item $x^2+2x-i\sqrt3 = 0$
\end{enumerate*}
\end{exercise}

\clearpage
\section{Vector spaces}

\begin{exercise}
Determine whether or not the following sets are vector spaces over the indicated fields:
\begin{enumerate}[label=(\alph*)]
\item $\Z$ over $\Z_2$, with addition of integers, and scaling defined as $0\cdot z = 0$, $1\cdot z = z$
for every $z\in\Z$
\item $\Z_4$ over $\Z_2$, with addition modulo $4$, and scaling defined as $0\cdot z = 0$,
$1\cdot z = z$ for every $z\in\Z_4$
\item $\C$ over $\R$, with addition of complex numbers, and scaling by real numbers
\item $\R$ over $\C$, with addition of real numbers, and scaling by complex numbers
\item all subsets of a finite set $X$ over $\Z_2$, with the symmetric difference as vector addition, and
with scaling defined as $0\cdot A = \emptyset$ and $1\cdot A = A$ for every $A\subseteq X$
\item $\R_+$ over $\R$, with addition of real numbers, and scaling defined as $p\odot r = r^p$ for every
$r\in\R_+$ and $p\in\R$
\end{enumerate}
\end{exercise}

\begin{exercise}
Decide which of the following subsets are subspaces:
\begin{enumerate}[label=(\alph*)]
\item $\{(x,y)\in\R^2 \mid x\geq 0,\ y\geq 0\}\subseteq \R^2$
\item $\{(x,y,z)\in\R^3 \mid yz \leq 0\}\subseteq \R^3$
\item $\{(x,y,z)\in\R^3 \mid x+y+z = x-y = 0\}\subseteq \R^3$
\item $\{(x,y)\in\R^2 \mid |x-y|\leq 1\}\subseteq \R^2$
\end{enumerate}
\end{exercise}

\begin{exercise}[$\ast$]
Prove that in a vector space $\V$ over a field $\K$, for every $\vv,\uu\in \V$ and $\lambda\in\K$,
\[
\lambda\cdot(\vv-\uu) = \lambda\cdot\vv - \lambda\cdot\uu.
\]
\end{exercise}

\begin{exercise}
The intersection of any collection of subspaces of a vector space $\V$ is a subspace of $\V$. Show that
this is not necessarily true for the union of subspaces.
\end{exercise}

\begin{exercise}
The real plane $\R^2$ is a vector space over $\R$. Describe, in geometrical terms, all subspaces of
$\R^2$.
\end{exercise}

\begin{exercise}[$\ast$]
Let $\V$ be a vector space over a field $\K$ and let $S,T\subseteq \V$. Prove that:
\begin{enumerate}[label=(\alph*)]
\item $S\subseteq \Span{S}$
\item $S\subseteq T \implies \Span{S}\subseteq\Span{T}$
\item $\Span{\Span{S}} = \Span{S}$
\item $\vv\in\Span{S} \iff \Span{S} = \Span{S\cup\{\vv\}}$.
\end{enumerate}
\end{exercise}

\clearpage
\section{Linear independence. Bases and dimension}

\begin{exercise}
Decide whether the given sequence of vectors is linearly independent in the indicated vector space:
\begin{enumerate}[label=(\alph*)]
\item $\big((1,0,1),(1,1,0),(0,1,1)\big)$ in $\R^3$ over $\R$
\item $\big((1,0,1),(1,1,0),(0,1,1)\big)$ in $\Z_2^3$ over $\Z_2$
\end{enumerate}
\end{exercise}

\begin{exercise}
Let $\V$ be any vector space over a field $\K$ and let $\vv_1,\vv_2,\ldots,\vv_n,\uu \in \V$. Decide
whether the given sequence of vectors is linearly independent:
\begin{enumerate}[label=(\alph*)]
\item $(\vv_1,\ \vv_1+\vv_2,\ \vv_1+\vv_2+\vv_3,\ \ldots,\ \vv_1+\ldots+\vv_{n-1}+\vv_n)$ if the sequence
$(\vv_1,\vv_2,\ldots,\vv_n)$ is linearly independent
\item $(\vv_1+\vv_2,\ \vv_2+\vv_3,\ \ldots,\ \vv_{n-1}+\vv_n,\ \vv_n+\vv_1)$ if the sequence
$(\vv_1,\vv_2,\ldots,\vv_n)$ is linearly independent
\item $(\vv_1+\uu,\ \vv_2+\uu,\ \ldots,\ \vv_n+\uu)$ if the sequence $(\vv_1,\vv_2,\ldots,\vv_n)$ is
linearly independent and $\uu\in \V$ is any vector
\item $S$, where $S$ is a subsequence of a linearly independent sequence $T$ of vectors from $\V$
\item $S$, where $S$ is a sequence of vectors from $\V$ containing a linearly independent subsequence $T$
\item ($\ast$) $(\vv_1,\vv_2,\ldots,\vv_n)$, where $\vv_i\in\Span{S_i}$ and $\vv_i\neq \0$ for
$i = 1,2,\ldots,n$, the sets $S_i\subseteq \V$ are finite and pairwise disjoint, and the set
$S_1\cup S_2\cup\ldots\cup S_n$ consists of linearly independent vectors.
\end{enumerate}
\end{exercise}

\begin{exercise}
Find the dimension of each of the following vector spaces:
\begin{enumerate}[label=(\alph*)]
\item $\C$ over $\R$
\item $\K^n$ over any field $\K$
\item $\Span{\{\vv_1,\vv_2,\ldots,\vv_{n-1}\}}$, where $(\vv_1,\vv_2,\ldots,\vv_n)$ is a basis of a
vector space $\V$ over a field $\K$
\item $\Span{\{\vv_n,\vv_{n-1},\ldots,\vv_1\}}$, where $(\vv_1,\vv_2,\ldots,\vv_n)$ is a basis of a
vector space $\V$ over a field $\K$
\item $\{(x,y,z,t)\in\R^4 \mid x+y+z+t = 0\}$
\item $\{(x,y,z,t)\in\R^4 \mid 2x+y = z-t\}$
\item $\Span{\{(1,2,1,0),(1,3,1,1),(0,1,2,-1),(-3,1,2,2)\}}$ over $\R$
\item $2^{\{a,b,c\}}$ over $\Z_2$ (with the symmetric difference as vector addition).
\end{enumerate}
\end{exercise}

\begin{exercise}
Calculate the coordinates of the vector $\vv$ with respect to the basis $S$:
\begin{enumerate}[label=(\alph*)]
\item $\vv = (1,2,3,4)$, $\ S = \big((1,1,1,1),(1,1,1,0),(1,1,0,1),(1,0,1,1)\big)$
\item $\vv = \{a,b\}$, $\ S = \big(\{a\},\{b\},\{c\}\big)$ in the vector space $2^{\{a,b,c\}}$ over
$\Z_2$
\item $\vv = \vv_1$, $\ S = (\vv_1,\vv_2,\ldots,\vv_n)$
\item $\vv = \vv_1+\vv_2+\ldots+\vv_n$, $\ S = (\vv_1,\vv_2,\ldots,\vv_n)$.
\end{enumerate}
\end{exercise}

\begin{exercise}[$\ast$]
Prove that if $(\vv_1,\vv_2,\ldots,\vv_n)$ is a basis of a vector space $\V$, then so is
$(\vv_1,\ \vv_1+\vv_2,\ \ldots,\ \vv_1+\vv_2+\ldots+\vv_n)$.
\end{exercise}

\begin{exercise}[$\ast$]
Prove that $\dim\Span{A\cup B} \leq \dim\Span{A} + \dim\Span{B}$ for subsets $A,B$ of a vector space
$\V$.
\end{exercise}

\clearpage
\section{Matrices}

\begin{exercise}
For the following matrices $A$ and $B$ calculate the matrix $A\cdot B$:
\begin{enumerate}[label=(\alph*)]
\item $A = \left( \begin{array}{cc} 1 & 0\\ 2 & 3\\ 1 & 2\\ -1 & 1 \end{array} \right)$, \quad
$B = \left( \begin{array}{ccc} 3 & -2 & 2\\ 1 & 5 & 3 \end{array} \right)$ \quad over the field $\R$
\item $A = \left( \begin{array}{ccc} 1 & 1 & 1\\ 1 & 1 & 1\\ 1 & 1 & 1 \end{array} \right)$, \quad
$B = \left( \begin{array}{ccc} 1 & 4 & 0\\ 2 & 5 & 1\\ 3 & 6 & 2 \end{array} \right)$ \quad over the
field $\Z_7$
\item $A = \left( \begin{array}{cccc} 1 & 1 & 0 & 0\\ 3 & 2 & 3 & 1\\ 1 & 0 & 1 & 0 \end{array} \right)$,
\quad $B = \left( \begin{array}{cc} 2 & 1\\ 1 & 0\\ 2 & 0\\ 1 & 1 \end{array} \right)$ \quad over the
field $\Z_5$
\end{enumerate}
\end{exercise}

\begin{exercise}
Recall that a square matrix $A$ is symmetric if and only if $A = A^T$. Is it true that multiplication of
symmetric matrices is commutative?
\end{exercise}

\begin{exercise}
Show for matrices $A$ and $B$ (of appropriate sizes) that $(A\cdot B)^T = B^T\cdot A^T$.
\end{exercise}

\begin{exercise}
Use matrices to determine the dimension of the subspace $W$ over the field $\R$:
\begin{enumerate}[label=(\alph*)]
\item $W = \Span{\{(1,2,3),(1,-1,2),(-1,4,-1)\}}$
\item $W = \Span{\{(3,1,2,1),(1,-1,1,2),(2,3,0,-3),(2,-1,1,2)\}}$
\end{enumerate}
\end{exercise}

\begin{exercise}
Calculate the following determinants:
\begin{enumerate}[label=(\alph*)]
\item $\det \left( \begin{array}{cccc}
1 & 2 & -2 & 1\\ -3 & 1 & 3 & 2\\ -1 & 2 & 4 & 2\\ 5 & -1 & 2 & 2
\end{array} \right)$, where the matrix belongs to $\R_4^4$
\item $\det \left( \begin{array}{cccc}
1 & 2 & 1 & 3\\ 2 & 1 & 2 & 1\\ 3 & 2 & 3 & 2\\ 4 & 3 & 4 & 3
\end{array} \right)$, where the matrix belongs to $(\Z_5)_4^4$
\item $\det(A)$, where $A\in \K_5^5$ has the entries $a_{ij} = i-j$.
\end{enumerate}
\end{exercise}

\begin{exercise}
The numbers $24026$, $40262$, $2624$, $26240$ and $62402$ are divisible by $41$. Prove that the following
determinant is also divisible by $41$:
\[
\det \left( \begin{array}{ccccc}
2 & 4 & 0 & 2 & 6\\
4 & 0 & 2 & 6 & 2\\
0 & 2 & 6 & 2 & 4\\
2 & 6 & 2 & 4 & 0\\
6 & 2 & 4 & 0 & 2
\end{array} \right).
\]
\end{exercise}

\begin{exercise}
Let $y\in\R$, $n\in\N$ and let $A\in\R_n^n$ be defined in the following way: $a_{ij} = y$ for
$j\neq n-i+1$ and $a_{ij} = 0$ for $j = n-i+1$. Find the determinant of $A$ for $n=3$, and then find the
general formula for the determinant of $A$ for an arbitrary $n\in\N$. Justify your answer.
\end{exercise}

\begin{exercise}
Show that for all matrices $A$ and $B$ of size $2\times 2$, $\ \det(A\cdot B) = \det(A)\cdot\det(B)$.
\end{exercise}

\begin{exercise}
Show, for invertible matrices $A$ and $B$, that $(A\cdot B)^{-1} = B^{-1}\cdot A^{-1}$.
\end{exercise}

\begin{exercise}
Calculate the inverse (if it exists) of each of the following matrices:
\[
A = \left( \begin{array}{cccc}
1 & 1 & 1 & 1\\ 0 & 1 & 1 & 1\\ 0 & 0 & 1 & 1\\ 0 & 0 & 0 & 1
\end{array} \right) \in (\Z_3)_4^4, \qquad
B = \left( \begin{array}{cccc}
1 & 2 & 2 & 1\\ 1 & 3 & 4 & 1\\ 0 & 3 & 1 & 0\\ 1 & 3 & 4 & 2
\end{array} \right) \in \R_4^4, \qquad
C = \left( \begin{array}{cccc}
1 & 2 & 2 & 1\\ 1 & 3 & 4 & 1\\ 0 & 3 & 1 & 0\\ 1 & 3 & 4 & 2
\end{array} \right) \in (\Z_5)_4^4.
\]
\end{exercise}

\clearpage
\section{Systems of linear equations}

\begin{exercise}
Find a basis and the dimension of the solution space over the field $\R$ of each of the following systems
of linear equations:
\[
\left\{ \begin{aligned}
x + 2y - 2z + 2s - t &= 0\\
x + 2y - z + 3s - 2t &= 0\\
2x + 4y - 7z + s + t &= 0
\end{aligned} \right.
\qquad\qquad
\left\{ \begin{aligned}
x + 2y - z + 3s - 4t &= 0\\
2x + 4y - 2z - s + 5t &= 0\\
2x + 4y - 2z + 4s - 2t &= 0
\end{aligned} \right.
\]
\end{exercise}

\begin{exercise}
Discuss the solvability of the following systems of equations over the field $\R$ in terms of the
parameters $a$ and $b$:
\[
\left\{ \begin{aligned}
ax + by + az + t &= 1\\
ax + y + az + t &= 1\\
x + y + az + at &= 1
\end{aligned} \right.
\qquad
\left\{ \begin{aligned}
x + 2y + 3z + 4t &= a\\
2x + 3y + 4z + 5t &= b\\
3x + 4y + 5z + 6t &= a\\
4x + 5y + 6z + 7t &= b
\end{aligned} \right.
\qquad
\left\{ \begin{aligned}
x + ay + az &= 1\\
ax + ay + z &= 1\\
ax + y + az &= 1\\
ax + ay + az &= 1
\end{aligned} \right.
\]
\end{exercise}

\begin{exercise}
Find general solutions of the following systems of equations over the field $\R$:
\begin{enumerate}[label=(\alph*)]
\item $\left\{ \begin{aligned}
2x + 2y + 3z + 4t &= 2\\
2x + 3y + 4z + 6t &= 0\\
3x + 4y + 5z + 6t &= 1\\
x + y + z + t &= 1
\end{aligned} \right.$
\item $\left\{ \begin{aligned}
x + y + z - t + u &= 1\\
x - y + z + t + u &= 1\\
-x + y - z + t - u &= 0
\end{aligned} \right.$
\item $\left\{ \begin{aligned}
x + y + z + t + u &= 0\\
x - 2y + 3z - 4t + 5u &= 0
\end{aligned} \right.$
\end{enumerate}
\end{exercise}

\begin{exercise}
Find all solutions in $\Z_7$ of the following system of equations:
\[
\left\{ \begin{aligned}
3x_1 + 4x_2 + 4x_3 + 6x_4 &= 1\\
4x_1 + 5x_2 + 4x_3 + x_4 &= 0\\
5x_1 + 6x_2 + 4x_3 + 3x_4 &= 6
\end{aligned} \right.
\qquad \text{over the field } \Z_7.
\]
\end{exercise}

\clearpage
\section{Linear mappings. Matrices of linear mappings}

\begin{exercise}
Verify which of the following mappings are linear:
\begin{enumerate}[label=(\alph*)]
\item $F\colon\R^3\to\R^2$, $\ F(x,y,z) = (x,\,z)$
\item $F\colon\R^2\to\R^2$, $\ F(x,y) = (3x+y+1,\ 2x-3y)$
\item $F\colon\R^2\to\R^3$, $\ F(x,y) = (x+2y,\ x-y,\ x)$
\item $F\colon\R\to\R$, $\ F(x) = |x|$
\item $F\colon\R^3\to\R^3$, $\ F(x,y,z) = (xy,\ x,\ z)$
\end{enumerate}
\end{exercise}

\begin{exercise}
Let $T\colon\C\to\C$, $\ T(z) = \bar z$.
\begin{enumerate}[label=(\alph*)]
\item Show that $T$ is \emph{not} linear if we consider $\C$ as a vector space over itself.
\item Show that $T$ is linear if we consider $\C$ as a vector space over $\R$.
\end{enumerate}
\end{exercise}

\begin{exercise}\label{ex:ker-im}
Find $\ker F$, $\operatorname{im} F$, $\rank(F)$ and the nullity of $F$ for the following linear mappings
$F\colon U\to V$:
\begin{enumerate}[label=(\alph*)]
\item $U = V = \R^3$, $\ F(x,y,z) = (x+y+z,\ x+y,\ x)$
\item $U = V = \R^3$, $\ F(x,y,z) = (x+y,\ x+y,\ x+2y-z)$
\item $U = \R^3$, $V = \R^2$, $\ F(x,y,z) = (x-y+z,\ x+y-z)$
\end{enumerate}
\end{exercise}

\begin{exercise}
Find the matrices of the linear mappings from Exercise~\ref{ex:ker-im} in the following bases of the
spaces $U$ and $V$, respectively:
\begin{enumerate}[label=(\alph*)]
\item the standard bases of unit vectors
\item $\big((1,1,1),(1,1,0),(1,0,0)\big)$ and $\big((1,0,0),(1,1,0),(1,1,1)\big)$
\item $\big((1,1,1),(1,0,0),(0,0,1)\big)$ and $\big((1,2),(3,3)\big)$
\end{enumerate}
\end{exercise}

\begin{exercise}
Find a formula and the matrix in the standard basis for the linear operator $F$ such that:
\begin{enumerate}[label=(\alph*)]
\item $F(1,1,1) = (1,0,1)$, $\ F(1,1,0) = F(0,1,1) = (1,1,1)$
\item $F(1,2,1) = (1,0,0)$, $\ F(1,0,1) = (0,1,0)$, $\ F(1,0,0) = (1,0,0)$
\item $F(1,0,0,0) = F(0,1,0,0) = F(0,0,1,0) = F(0,0,0,1) = (1,1,1,1)$
\end{enumerate}
\end{exercise}

\begin{exercise}\label{ex:matrix-bases}
Assuming that $A = M_S(F)$ is the matrix of a linear operator $F$ in the basis $S$, find the matrix
$B = M_R(F)$ of $F$ in the basis $R$:
\begin{enumerate}[label=(\alph*)]
\item $\left( \begin{array}{cc} 1 & 1\\ 1 & 2 \end{array} \right)$, for $S = \big((1,1),(1,2)\big)$ and
$R = \big((1,0),(0,1)\big)$
\item $\left( \begin{array}{ccc} 1 & 0 & 0\\ 0 & 1 & 0\\ 0 & 0 & 1 \end{array} \right)$, for
$S = \big((1,1,1),(0,1,2),(1,0,1)\big)$ and $R = \big((1,0,0),(0,1,0),(0,0,1)\big)$
\item $\left( \begin{array}{ccc} 1 & 2 & 2\\ 1 & 1 & 1\\ 2 & 0 & 0 \end{array} \right)$, for
$S = \big((1,1,1),(0,1,2),(1,0,1)\big)$ and $R = \big((1,0,0),(0,1,0),(0,0,1)\big)$
\item $\left( \begin{array}{cccc} 1 & 1 & 1 & 1\\ 0 & 1 & 1 & 1\\ 0 & 0 & 1 & 1\\ 0 & 0 & 0 & 1
\end{array} \right)$, for $S = \big((1,0,1,0),(0,0,1,1),(1,1,0,0),(1,0,1,1)\big)$ and\\
$R = \big((1,1,1,0),(1,1,0,1),(1,0,1,1),(0,1,1,1)\big)$
\end{enumerate}
\end{exercise}

\begin{exercise}
For each pair of matrices $A$ and $B$ from Exercise~\ref{ex:matrix-bases}, find a matrix $P$ such that
$A = P^{-1}\cdot B\cdot P$. Verify your solution by matrix multiplication.
\end{exercise}

\clearpage
\section{Eigenvalues and eigenvectors}

\begin{exercise}
Find the eigenvalues of $A = M_S(F)$:
\begin{enumerate}[label=(\alph*)]
\item $A = \left( \begin{array}{cc} 2+i & 1\\ 2 & 2-i \end{array} \right)$
\item $A = \left( \begin{array}{ccc} 2 & -1 & 2\\ 5 & -3 & 3\\ -1 & 0 & -2 \end{array} \right)$
\end{enumerate}
\end{exercise}

\begin{exercise}
Find all eigenvalues and a basis of each eigenspace for the following linear operators:
\begin{enumerate}[label=(\alph*)]
\item $F\colon\R^2\to\R^2$, $\ F(x,y) = (x,\ x+y)$
\item $F\colon\R^3\to\R^3$, $\ F(x,y,z) = (2x+y,\ y-z,\ 2y+4z)$
\item $F\colon\R^3\to\R^3$, $\ F(x,y,z) = (x-z,\ 2y,\ x+z)$
\item $F\colon\C^2\to\C^2$, $\ F(x,y) = (-y,\ x)$
\item $F\colon\C^3\to\C^3$, $\ F(x,y,z) = (x-z,\ 2y,\ x+z)$.
\end{enumerate}
\end{exercise}

\begin{exercise}\label{ex:eig-matrices}
Assuming that each of the following matrices represents a linear operator $F$ in a basis $S$, find all
eigenvectors and a basis of each eigenspace:
\begin{enumerate}[label=(\alph*)]
\item $A = M_S(F) = \left( \begin{array}{ccc} 3 & 1 & 1\\ 2 & 4 & 2\\ 1 & 3 & 3 \end{array} \right)$
\item $A = M_S(F) = \left( \begin{array}{ccc} 1 & 2 & 2\\ 1 & 2 & -1\\ -1 & 1 & 4 \end{array} \right)$
\item $A = M_S(F) = \left( \begin{array}{ccc} 1 & 1 & 0\\ 0 & 1 & 0\\ 0 & 0 & 1 \end{array} \right)$
\end{enumerate}
\end{exercise}

\begin{exercise}
For each operator $F$ from Exercise~\ref{ex:eig-matrices}, determine whether there exists a basis $T$
such that the matrix of the operator $F$ in the basis $T$ is diagonal.
\end{exercise}

\begin{exercise}
For the matrix
$A = \left( \begin{array}{ccc} 1 & 2 & -1\\ 1 & -1 & 0\\ 2 & -2 & 0 \end{array} \right)$
calculate $A^{100}$.
\end{exercise}

\begin{exercise}
For a linear mapping $F\colon\R^2\to\R^2$ such that $F(1,2) = (2,4)$ and $F(-2,1) = (-2,1)$, calculate
$F^{150}(1,0)$.
\end{exercise}
